\documentclass[a4paper,12pt]{article}
\usepackage[utf8]{inputenc}
\usepackage[margin=1in]{geometry}

\usepackage{amsmath}
\usepackage{graphicx}
\graphicspath{{figures/}}
\usepackage{tikz}
\usepackage{hyperref}
\hypersetup{
    colorlinks=true,
    linkcolor=blue,
    filecolor=magenta,      
    urlcolor=cyan,
    }
\urlstyle{same}

\renewcommand{\vec}[1]{\boldsymbol{#1}}
\newcommand{\unitvec}[1]{\hat{\vec{#1}}}
\newcommand{\mrm}[1]{\mathrm{#1}}
\newcommand{\diff}{\operatorname{d}\!}
\newcommand{\mat}[1]{\mathbf{#1}}
\newcommand{\unit}[1]{\ensuremath{\,\mrm{#1}}}
\renewcommand{\Re}{\operatorname{Re}}
\renewcommand{\Im}{\operatorname{Im}}

\newcommand{\iu}{\mrm{i}}
\newcommand{\ju}{\mrm{j}}
\newcommand{\eu}{\mrm{e}}
\newcommand{\cu}{\mrm{c}_{0}}
\newcommand{\Ev}{\vec{E}}
\newcommand{\Hv}{\vec{H}}
\newcommand{\Fv}{\vec{F}}
\newcommand{\Av}{\vec{A}}
\newcommand{\Kv}{\vec{K}}
\newcommand{\pv}{\vec{p}}
\newcommand{\rv}{\vec{r}}
\newcommand{\kv}{\vec{k}}
\newcommand{\av}{\vec{a}}
\newcommand{\Zv}{\vec{0}}
\newcommand{\xuv}{\unitvec{x}}
\newcommand{\yuv}{\unitvec{y}}
\newcommand{\zuv}{\unitvec{z}}
\newcommand{\ruv}{\unitvec{r}}
\newcommand{\nuv}{\unitvec{n}}
\newcommand{\tuv}{\unitvec{t}}
\newcommand{\kuv}{\unitvec{k}}
\newcommand{\uuv}{\unitvec{u}}
\newcommand{\rhouv}{\unitvec{\rho}}
\newcommand{\varphiuv}{\unitvec{\varphi}}
\newcommand{\puv}{\unitvec{p}}
\newcommand{\suv}{\unitvec{s}}
\newcommand{\IM}{\mat{I}}
\newcommand{\ZM}{\mat{Z}}
\newcommand{\SM}{\mat{S}}
\newcommand{\TM}{\mat{T}}
\newcommand{\rmat}{\mat{r}}

\newcommand{\thetauv}{\unitvec{\theta}}
\newcommand{\phiuv}{\unitvec{\phi}}

\newcommand{\BesselJ}{\operatorname{J}}

\newcommand{\uv}{\vec{u}}
\newcommand{\vv}{\vec{v}}
\newcommand{\epsm}{\boldsymbol{\epsilon}}
\newcommand{\mum}{\boldsymbol{\mu}}

\newcommand{\ordo}{\operatorname{o}}
\newcommand{\Ordo}{\operatorname{O}}

\usepackage{listings}

\usepackage{xcolor}

\definecolor{codegreen}{rgb}{0,0.6,0}
\definecolor{codegray}{rgb}{0.5,0.5,0.5}
\definecolor{codepurple}{rgb}{0.58,0,0.82}
\definecolor{backcolour}{rgb}{0.95,0.95,0.92}

\lstdefinestyle{mystyle}{
    backgroundcolor=\color{backcolour},   
    commentstyle=\color{codegreen},
    keywordstyle=\color{magenta},
    numberstyle=\tiny\color{codegray},
    stringstyle=\color{codepurple},
    basicstyle=\ttfamily\footnotesize,
    breakatwhitespace=false,         
    breaklines=true,                 
    captionpos=b,                    
    keepspaces=true,                 
    numbers=left,                    
    numbersep=5pt,                  
    showspaces=false,                
    showstringspaces=false,
    showtabs=false,                  
    tabsize=2
}

\lstset{style=mystyle}




\title{Spherical vector waves}

\author{Daniel Sj\"oberg}

\begin{document}
\maketitle

\begin{abstract}
  This document covers the details of a python implementation of
  spherical harmonics and spherical vector waves based on Hansen's
  formulation.
\end{abstract}


\section{Introduction}

\section{Definitions from Hansen \cite{Hansen1988}}

\begin{itemize}
\item Wave \textbf{ADMITTANCE} is denoted
  \begin{equation}
    \eta = \sqrt{\epsilon/\mu}
  \end{equation}
\item Time convention
  \begin{equation}
    \eu^{-\iu\omega t}
  \end{equation}
\item Generating function
  \begin{equation}
    F_{mn}^{(c)}(r,\theta,\phi) = \frac{1}{\sqrt{2\pi}} \frac{1}{\sqrt{n(n+1)}} \left(-\frac{m}{|m|}\right)^{m} z_{n}^{(c)}(kr) \bar{P}_{n}^{|m|}(\cos\theta) \eu^{\iu m\phi}
  \end{equation}
  with $(-m/|m|)^{m} = 1$ for $m=0$.
\item Normalized associated Legendre functions
  \begin{equation}
    \bar{P}_{n}^{m}(\cos\theta) = \sqrt{\frac{2n+1}{2}\frac{(n-m)!}{(n+m)!}} P_{n}^{m}(\cos\theta)
  \end{equation}
  with (note the absence of the Condon-Shortley phase factor $(-1)^{m}$)
  \begin{equation}
    P_{n}^{m}(\cos\theta) = (\sin\theta)^{m}\frac{\diff^{m}P_{n}(\cos\theta)}{\diff(\cos\theta)^{m}}
  \end{equation}
\item Derivative of associated Legendre function ($\frac{\diff P_{n}^{m}(\cos\theta)}{\diff\theta} = -\sin\theta\frac{\diff P_{n}^{m}(\cos\theta)}{\diff(\cos\theta)}$)
  % \begin{equation}
  %   \frac{\diff P_{n}^{m}(\cos\theta)}{\diff\theta} =
  %   \begin{cases}
  %     -P_{n}^{1}(\cos\theta) & m=0 \\
  %     \frac{1}{2}\{(n-m+1)(n+m)P_{n}^{m-1}(\cos\theta) -P_{n}^{m+1}(\cos\theta) & m>0
  %   \end{cases}
  % \end{equation}
  \begin{equation}
    \frac{\diff P_{n}^{m}(x)}{\diff x} =
    \begin{cases}
      -\frac{1}{\sqrt{1-x^{2}}}P_{n}^{1}(x) & m=0 \\
      \frac{(n + m)(n - m + 1)}{\sqrt{1 - x^{2}}} P_{n}^{m-1}(x) + \frac{mx}{1-x^{2}}P_{n}^{m}(x) & m >0
    \end{cases}
  \end{equation}
\item Radial functions
  \begin{equation}
    R_{sn}^{(c)}(kr) =
    \begin{cases}
      z_{n}^{(c)}(kr), & s=1 \\
      \frac{1}{kr}\frac{\diff}{\diff(kr)}\{krz_{n}^{(c)}(kr)\}, & s=2
    \end{cases}
  \end{equation}
  where
  \begin{align}
    z_{n}^{(1)}(kr) &= j_{n}(kr) && \text{(spherical Bessel function)} \\
    z_{n}^{(2)}(kr) &= y_{n}(kr) && \text{(spherical Neumann function)} \\
    z_{n}^{(3)}(kr) &= h_{n}^{(1)}(kr) = j_{n}(kr) + \iu y_{n}(kr) && \text{(spherical Hankel function of the first kind)} \\
    z_{n}^{(4)}(kr) &= h_{n}^{(2)}(kr) = j_{n}(kr) - \iu y_{n}(kr) && \text{(spherical Hankel function of the second kind)}     
  \end{align}
\item Spherical vector waves
  \begin{multline}
    \Fv^{(c)}_{1mn}(r,\theta,\phi) = \nabla F_{mn}^{(c)}(r,\theta,\phi)\times\rv \\
    = \frac{1}{\sqrt{2\pi}} \frac{1}{\sqrt{n(n+1)}} \left(-\frac{m}{|m|}\right)^{m} \bigg\{ z_{n}^{(c)}(kr) \frac{\iu m\bar{P}_{n}^{|m|}(\cos\theta)}{\sin\theta} \eu^{\iu m\phi}\thetauv \\
    - z_{n}^{(c)}(kr) \frac{\diff \bar{P}_{n}^{|m|}(\cos\theta)}{\diff\theta} \eu^{\iu m\phi}\phiuv \bigg\}
  \end{multline}
  \begin{multline}
    \Fv^{(c)}_{2mn}(r,\theta,\phi) = \frac{1}{k}\nabla\times\Fv_{1mn}^{(c)}(r,\theta,\phi) \\
    = \frac{1}{\sqrt{2\pi}} \frac{1}{\sqrt{n(n+1)}} \left(-\frac{m}{|m|}\right)^{m} \bigg\{ \frac{n(n+1)}{kr} z_{n}^{(c)}(kr) \bar{P}_{n}^{|m|}(\cos\theta) \eu^{\iu m\phi}\ruv \\
    + \frac{1}{kr}\frac{\diff}{\diff(kr)} \left\{krz_{n}^{(c)}(kr)\right\} \frac{\diff \bar{P}_{n}^{|m|}(\cos\theta)}{\diff\theta} \eu^{\iu m\phi}\thetauv \\
    + \frac{1}{kr}\frac{\diff}{\diff(kr)} \left\{krz_{n}^{(c)}(kr)\right\} \frac{\iu m\bar{P}_{n}^{|m|}(\cos\theta)}{\sin\theta} \eu^{\iu m\phi} \phiuv \bigg\}
  \end{multline}
  (A third spherical vector wave
  $\Fv_{3mn}^{(c)}(r,\theta,\phi) \sim \nabla
  F_{mn}^{(c)}(r,\theta,\phi)$ can be defined, but is not used in
  \cite{Hansen1988}.)
\item Conjugation relations (assuming $kr$ is real)
  \begin{align}
    \Fv_{smn}^{(1)*}(r,\theta,\phi) &= (-1)^{m}\Fv_{s,-m,n}^{(1)}(r,\theta,\phi) \\
    \Fv_{smn}^{(2)*}(r,\theta,\phi) &= (-1)^{m}\Fv_{s,-m,n}^{(2)}(r,\theta,\phi) \\
    \Fv_{smn}^{(3)*}(r,\theta,\phi) &= (-1)^{m}\Fv_{s,-m,n}^{(4)}(r,\theta,\phi) \\
    \Fv_{smn}^{(4)*}(r,\theta,\phi) &= (-1)^{m}\Fv_{s,-m,n}^{(3)}(r,\theta,\phi)
  \end{align}
  Summary:
  $\Fv_{smn}^{(c)*}(r,\theta,\phi) =
  (-1)^{m}\Fv_{s,-m,n}^{(c')}(r,\theta,\phi)$, where $c'=1,2$ if
  $c=1,2$, and $c'=4,3$ if $c=3,4$.
\item Expansion of fields and power
  \begin{align}
    \Ev(r,\theta,\phi) &= \frac{k}{\sqrt{\eta}}\sum_{csmn} Q_{smn}^{(c)}\Fv_{smn}^{(c)}(r,\theta,\phi) \\
    \Hv(r,\theta,\phi) &= -\iu k\sqrt{\eta}\sum_{csmn} Q_{smn}^{(c)}\Fv_{3-s,m,n}^{(c)}(r,\theta,\phi) \\
    P &= \oint\frac{1}{2}\Re(\Ev\times\Hv^{*})\cdot\nuv\diff S = \frac{1}{2}\sum_{smn}|Q_{smn}^{(3)}|^{2}
  \end{align}
  \begin{multline}
    P = \oint\frac{1}{2}\Re(\Ev\times\Hv^{*})\cdot\nuv\diff S = \frac{1}{4} \oint (\Ev\times\Hv^{*} + \Ev^{*}\times\Hv) \cdot \nuv \diff S\\
    = \frac{1}{4}\sum_{csmn}\sum_{\gamma\sigma\mu\nu} \oint \left\{ Q_{smn}^{(c)}\Ev_{smn}^{(c)}\times(Q_{\sigma\mu\nu}^{(\gamma)}\Hv_{\sigma\mu\nu}^{(\gamma)})^{*} + (Q_{\sigma\mu\nu}^{(\gamma)}\Ev_{\sigma\mu\nu}^{(\gamma)})^{*}\times Q_{smn}^{(c)}\Hv_{smn}^{(c)} \right\} \cdot\nuv\diff S 
  \end{multline}
  Since
  $\Ev_{\sigma\mu\nu}^{(\gamma)*} =
  (-1)^{\mu}\frac{k}{\sqrt{\eta}}\Fv_{\sigma,-\mu,\nu}^{(\gamma')} =
  (-1)^{\mu}\Ev_{\sigma,-\mu,\nu}^{(\gamma')}$ and
  $\Hv_{\sigma\mu\nu}^{(\gamma)*} = (-1)^{\mu}\iu
  k\sqrt{\eta}\Fv_{\sigma,-\mu,\nu}^{(\gamma')} =
  -(-1)^{\mu}\Hv_{\sigma,-\mu,\nu}^{(\gamma')}$, we have
  \begin{multline}
    P = \frac{1}{4} \sum_{csmn}\sum_{\gamma\sigma\mu\nu} Q_{smn}^{(c)}Q_{\sigma\mu\nu}^{(\gamma)*} (-1)^{\mu} \oint \left\{ -\Ev_{smn}^{(c)}\times\Hv_{\sigma,-\mu,\nu}^{(\gamma')} + \Ev_{\sigma,-\mu,\nu}^{(\gamma')}\times\Hv_{smn}^{(c)} \right\} \cdot \nuv \diff S \\
    = \frac{1}{4} \sum_{csmn}\sum_{\gamma\sigma\mu\nu} Q_{smn}^{(c)}Q_{\sigma\mu\nu}^{(\gamma)*} (-1)^{\mu} (-\delta_{s\sigma}\delta_{m,-(-\mu)}\delta_{n\nu}(-1)^{m}(-\iu)A^{(c,\gamma')}) \\
    = \frac{1}{4} \sum_{c\gamma smn} Q_{smn}^{(c)}Q_{smn}^{(\gamma)*} \iu A^{(c,\gamma')}
    = \sum_{smn}\frac{\iu}{4} [ Q_{smn}^{(1)}Q_{smn}^{(2)*} + Q_{smn}^{(1)}Q_{smn}^{(3)*}(-\iu) + Q_{smn}^{(1)}Q_{smn}^{(4)*}\iu \\
    + Q_{smn}^{(2)}Q_{smn}^{(1)*}(-1) + Q_{smn}^{(2)}Q_{smn}^{(3)*}(-1) + Q_{smn}^{(2)}Q_{smn}^{(4)*}(-1) \\
    + Q_{smn}^{(3)}Q_{smn}^{(1)*}(-\iu) + Q_{smn}^{(3)}Q_{smn}^{(2)*} + Q_{smn}^{(3)}Q_{smn}^{(3)*}(-2\iu) \\
    + Q_{smn}^{(4)}Q_{smn}^{(1)*}\iu + Q_{smn}^{(4)}Q_{smn}^{(2)*} + Q_{smn}^{(4)}Q_{smn}^{(4)*}(2\iu) ] \\
    = \frac{1}{2} \sum_{smn} \Big[ -\Im(Q_{smn}^{(1)}Q_{smn}^{(2)*}) + \Re(Q_{smn}^{(1)}Q_{smn}^{(3)*}) - \Re(Q_{smn}^{(1)}Q_{smn}^{(4)*}) \\
    + \Im(Q_{smn}^{(2)}Q_{smn}^{(3)*}) + \Im(Q_{smn}^{(2)}Q_{smn}^{(4)*}) + |Q_{smn}^{(3)}|^{2} - |Q_{smn}^{(4)}|^{2} \Big] \\
    = \frac{1}{2} \sum_{smn} \Big[ -\Im(Q_{smn}^{(1)}Q_{smn}^{(2)*}) + \Re(Q_{smn}^{(1)}(Q_{smn}^{(3)*} - Q_{smn}^{(4)*})) + \Im(Q_{smn}^{(2)}(Q_{smn}^{(3)*} + Q_{smn}^{(4)*})) \\
      + |Q_{smn}^{(3)}|^{2} - |Q_{smn}^{(4)}|^{2} \Big]
  \end{multline}
  which is
  $P = \frac{1}{2}\sum_{smn}|Q_{smn}^{(3)}|^{2} - |Q_{smn}^{(4)}|^{2}$
  if $Q_{smn}^{(1)} = Q_{smn}^{(2)} = 0$.
\item Single index convention
  \begin{equation}
    j = 2\{n(n+1) + m-1\} + s, \quad j = 1, 2, \ldots, J, \quad\text{where $J = 2N(N+2)$}
  \end{equation}
  with inverse
  \begin{align}
    s &= \begin{cases} 1 & \text{$j$ odd} \\ 2 & \text{$j$ even} \end{cases} \\
    n &= \text{integer part of $\displaystyle\sqrt{\frac{j-s}{2}+1}$} \\
    m &= \frac{j-s}{2} + 1 - n(n+1) 
  \end{align}
\item Representation of far fields
  \begin{align}
    \Kv_{smn}(\theta,\phi) &= \lim_{kr\rightarrow\infty} \left\{ \sqrt{4\pi}\frac{kr}{\eu^{\iu kr}} \Fv_{smn}^{(3)}(r,\theta,\phi) \right\} \\
    \Kv_{1mn}(\theta,\phi) &= \sqrt{\frac{2}{n(n+1)}} \left(-\frac{m}{|m|}\right)^{m} \eu^{\iu m\phi} (-\iu)^{n+1} \left\{ \frac{\iu m\bar{P}_{n}^{|m|}(\cos\theta)}{\sin\theta}\thetauv - \frac{\diff\bar{P}_{n}^{|m|}(\cos\theta)}{\diff\theta}\phiuv \right\} \\
    \Kv_{2mn}(\theta,\phi) &= \sqrt{\frac{2}{n(n+1)}} \left(-\frac{m}{|m|}\right)^{m} \eu^{\iu m\phi} (-\iu)^{n} \left\{ \frac{\diff\bar{P}_{n}^{|m|}(\cos\theta)}{\diff\theta}\thetauv + \frac{\iu m\bar{P}_{n}^{|m|}(\cos\theta)}{\sin\theta}\phiuv \right\} 
  \end{align}
\item Orthogonality

  Product of radial components
  \begin{multline}
    \int_{\phi=0}^{2\pi}\int_{\theta=0}^{\pi} \{\Fv_{smn}^{(c)}(r,\theta,\phi)\cdot\ruv\} \{\Fv_{\sigma\mu\nu}^{(\gamma)}(r,\theta,\phi)\cdot\ruv\} \sin\theta\diff\theta\diff\phi \\
    = \delta_{s\sigma}\delta_{s2}\delta_{m,-\mu}\delta_{n\nu}(-1)^{m}n(n+1)\frac{z_{n}^{(c)}(kr)}{kr}\frac{z_{n}^{(\gamma)}(kr)}{kr}
  \end{multline}
  Product of tangential components
  \begin{multline}
    \int_{\phi=0}^{2\pi}\int_{\theta=0}^{\pi} \Big[ \{\Fv_{smn}^{(c)}(r,\theta,\phi)\cdot\thetauv\} \{\Fv_{\sigma\mu\nu}^{(\gamma)}(r,\theta,\phi)\cdot\thetauv\} \\
    + \{\Fv_{smn}^{(c)}(r,\theta,\phi)\cdot\phiuv\} \{\Fv_{\sigma\mu\nu}^{(\gamma)}(r,\theta,\phi)\cdot\phiuv\} \Big] \sin\theta\diff\theta\diff\phi \\
    = \delta_{s\sigma}\delta_{m,-\mu}\delta_{n\nu}(-1)^{m}R_{sn}^{(c)}(kr)R_{sn}^{(\gamma)}(kr)
    \label{eq:orthotang}
  \end{multline}
  Scalar product
  \begin{multline}
    \int_{\phi=0}^{2\pi}\int_{\theta=0}^{\pi} \Fv_{smn}^{(c)}(r,\theta,\phi) \cdot \Fv_{\sigma\mu\nu}^{(\gamma)}(r,\theta,\phi) \sin\theta\diff\theta\diff\phi \\
    = \delta_{s\sigma}\delta_{m,-\mu}\delta_{n\nu}(-1)^{m} \left\{ R_{sn}^{(c)}(kr)R_{sn}^{(\gamma)}(kr) + \delta_{s2}n(n+1)\frac{z_{n}^{(c)}(kr)}{kr}\frac{z_{n}^{(\gamma)}(kr)}{kr} \right\}
  \end{multline}
  Vector product  
  \begin{multline}
    \int_{\phi=0}^{2\pi}\int_{\theta=0}^{\pi} \{ \Fv_{smn}^{(c)}(r,\theta,\phi) \times \Fv_{\sigma\mu\nu}^{(\gamma)}(r,\theta,\phi) \} \cdot \ruv \sin\theta\diff\theta\diff\phi \\
    = -\delta_{s,3-\sigma}\delta_{m,-\mu}\delta_{n\nu}(-1)^{m+s} R_{sn}^{(c)}(kr)R_{3-s,n}^{(\gamma)}(kr)
  \end{multline}
\item Reciprocity integral

  For the two fields, both representing a single mode with wave
  coefficients equal to 1 watts$^{1/2}$
  \begin{align}
    (\Ev_{smn}^{(c)},\Hv_{smn}^{(c)}) &= \left( \frac{k}{\sqrt{\eta}}\Fv_{smn}^{(c)}(r,\theta,\phi), -\iu k\sqrt{\eta}\Fv_{3-s,m,n}^{(c)}(r,\theta,\phi) \right) \\
    (\Ev_{\sigma\mu\nu}^{(\gamma)},\Hv_{\sigma\mu\nu}^{(\gamma)}) &= \left( \frac{k}{\sqrt{\eta}}\Fv_{\sigma\mu\nu}^{(\gamma)}(r,\theta,\phi), -\iu k\sqrt{\eta}\Fv_{3-\sigma,\mu,\nu}^{(\gamma)}(r,\theta,\phi) \right) 
  \end{align}
  we have
  \begin{equation}
    \oint_{S} \left\{ \Ev_{smn}^{(c)}\times\Hv_{\sigma\mu\nu}^{(\gamma)} - \Ev_{\sigma\mu\nu}^{(\gamma)}\times\Hv_{smn}^{(c)} \right\} \cdot \nuv\diff S = \delta_{s\sigma}\delta_{m,-\mu}\delta_{n\nu}(-1)^{m}(-\iu)A^{(c,\gamma)}
  \end{equation}
  where
  \begin{center}
    \begin{tabular}{c|cccc}
      $A^{(c,\gamma)}$ & $\gamma=1$ & $\gamma=2$ & $\gamma=3$ & $\gamma=4$ \\ \hline
      $c=1$ & $0$ & $1$ & $\iu$ & $-\iu$ \\
      $c=2$ & $-1$ & $0$ & $-1$ & $-1$ \\
      $c=3$ & $-\iu$ & $1$ & $0$ & $-2\iu$ \\
      $c=4$ & $\iu$ & $1$ & $2\iu$ & $0$ 
    \end{tabular}
  \end{center}
  Due to the reciprocity theorem, the surface $S$ can be deformed into
  another surface $S'$ without changing the value of the integral, as
  long as there are no sources between $S$ and $S'$.

  Hence, if we know the outgoing field
  $(\Ev(\rv),\Hv(\rv)) =
  \sum_{smn}Q_{smn}^{(3)}(\Ev_{smn}^{(3)},\Hv_{smn}^{(3)})$ on a closed
    surface $S$, we can compute the expansion coefficients from the
    relation
  \begin{equation}
    \oint_{S} \left\{ \Ev\times\Hv_{s,-m,n}^{(1)} - \Ev_{s,-m,n}^{(1)}\times\Hv \right\} \cdot \nuv\diff S = (-1)^{m+1}Q_{smn}^{(3)}
  \end{equation}

\item Expansion coefficients of a plane wave (valid for $r<R\approx N/k$)
  \begin{align}
    \kv_{0} &= -k\sin\theta_{0}\cos\phi_{0}\xuv - k\sin\theta_{0}\sin\phi_{0}\yuv - k\cos\theta_{0}\zuv \\
    \Ev_{0}\eu^{\iu\kv_{0}\cdot\rv} &= \frac{k}{\sqrt{\eta}} \sum_{s=1}^{2}\sum_{n=1}^{N}\sum_{m=-n}^{n} Q_{smn}\Fv_{smn}^{(1)}(r,\theta,\phi) \\
    Q_{smn} &= \frac{\sqrt{\eta}}{k} (-1)^{m}\sqrt{4\pi}\iu\Ev_{0}\cdot\Kv_{s,-m,n}(\theta_{0},\phi_{0})
  \end{align}

\item Rotation: The original coordinate system is $(x,y,z)$, the final
  is $(x',y',z')$. Three consecutive rotations (Euler angles) are
  described as follows, where $(x_{1},y_{1},z_{1})$ is the coordinate
  system after the first rotation, and $(x_{2},y_{2},z_{2})$ is the
  coordinate system after the second rotation.
  \begin{enumerate}
  \item A rotation about the $z$-axis (through an angle $\phi_{0}$).
  \item A rotation about the $y_{1}$-axis (through an angle $\theta_{0}$).
  \item A rotation about the $z_{2}$-axis (through an angle $\chi_{0}$). 
  \end{enumerate}
  The spherical vector waves are rotated as
  \begin{equation}
    \Fv_{smn}^{(c)}(r,\theta,\phi) = \sum_{\mu=-n}^{n} \eu^{\iu m\phi_{0}} d_{\mu m}^{n}(\theta_{0})\eu^{\iu\mu\chi_{0}} \Fv_{s\mu n}^{(c)}(r',\theta',\phi')
  \end{equation}
  where (the symbol $\binom{i}{j} = \frac{i!}{(i-j)!j!}$ is the
  binomial coefficient and the summation over $\sigma$ is such that no
  factorial have a negative argument)
  \begin{multline}
    d_{\mu m}^{n}(\theta) = \left\{\frac{(n+\mu)!(n-\mu)!}{(n+m)!(n-m)!}\right\}^{1/2} \sum_{\sigma} \binom{n+m}{n-\mu-\sigma}\binom{n-m}{\sigma} (-1)^{n-\mu-\sigma} \\
    \cdot\left(\cos\frac{\theta}{2}\right)^{2\sigma+\mu+m} \left(\sin\frac{\theta}{2}\right)^{2n-2\sigma-\mu-m}
  \end{multline}
  The rotation coefficient satisfies the normalization property
  \begin{equation}
    \sum_{\mu=-n}^{n}(d_{\mu m}^{n}(\theta))^{2} = 1, \quad \text{all $(m,n)$}
  \end{equation}
  and the orthogonality relation
  \begin{equation}
    \int_{0}^{\pi} d_{\mu m}^{n}(\theta)d_{\mu m}^{n'}(\theta)\sin\theta\diff\theta = \frac{2}{2n+1}\delta_{nn'}
  \end{equation}
  The transformation of expansion coefficients is given by
  \begin{multline}
    \Ev(\rv) = \sum_{csmn} Q_{smn}^{(c)}\Fv_{smn}^{(c)}(r,\theta,\phi) = \sum_{csmn} Q_{smn}^{(c)}\sum_{\mu=-n}^{n}\eu^{\iu m\phi_{0}} d_{\mu m}^{n}(\theta_{0})\eu^{\iu\mu\chi_{0}} \Fv_{s\mu n}^{(c)}(r',\theta',\phi') \\
    = \sum_{cs\mu n}\underbrace{\left[\sum_{m} Q_{smn}^{(c)}\eu^{\iu m\phi_{0}} d_{\mu m}^{n}(\theta_{0})\eu^{\iu\mu\chi_{0}}\right]}_{=Q_{s\mu n}^{(c)'}} \Fv_{s\mu n}^{(c)}(r',\theta',\phi') 
  \end{multline}
  
\item Translation: A spherical vector wave function defined in the
  unprimed coordinate system $(r,\theta,\phi)$ may be expressed as a
  combination of spherical waves defined in the primed system
  $(r',\theta',\phi')$ (which is translated a distance $A$ in the
  positive direction of the $z$-axis)
  \begin{align}
    \Fv_{s\mu n}^{(c)}(r,\theta,\phi) &= \sum_{\sigma=1}^{2} \sum_{\nu=|\mu|,\nu\neq0}^{\infty} C_{\sigma\mu\nu}^{sn(c)}(kA) \Fv_{\sigma\mu\nu}^{(1)}(r',\theta',\phi') & r'&<|A| \\
    \Fv_{s\mu n}^{(c)}(r,\theta,\phi) &= \sum_{\sigma=1}^{2} \sum_{\nu=|\mu|,\nu\neq0}^{\infty} C_{\sigma\mu\nu}^{sn(1)}(kA) \Fv_{\sigma\mu\nu}^{(c)}(r',\theta',\phi') & r'&>|A| 
  \end{align}
  where
  \begin{multline}
    C_{\sigma\mu\nu}^{sn(c)}(kA) = \sqrt{\frac{(2n+1)(2\nu+1)}{n(n+1)\nu(\nu+1)}} \sqrt{\frac{(\nu+\mu)!(n-\mu)!}{(\nu-\mu)!(n+\mu)!}} (-1)^{\mu} \frac{1}{2} \iu^{n-\nu} \\
    \cdot \sum_{p=|n-\nu|}^{n+\nu} \Big[ \iu^{-p}(\delta_{s\sigma}\{n(n+1) + \nu(\nu+1) - p(p+1)\} + \delta_{3-s,\sigma}\{2\iu\mu kA\}) \\
    \cdot a(\mu,n,-\mu,\nu,p) z_{p}^{(c)}(kA) \Big]
  \end{multline}
  and the linearization coefficients are
  \begin{equation}
    a(\mu,n,-\mu,\nu,p) = (2p+1) \left\{\frac{(n+\mu)!(\nu-\mu)!}{(n-\mu)!(\nu+\mu)!}\right\}^{1/2}
    \begin{pmatrix}
      n & \nu & p \\
      0 & 0 & 0
    \end{pmatrix}
    \begin{pmatrix}
      n & \nu & p \\
      \mu & -\mu & 0
    \end{pmatrix}
  \end{equation}
  where the Wigner 3-j symbols are
  \begin{multline}
    \begin{pmatrix}
      j_{1} & j_{2} & j_{3} \\
      m_{1} & m_{2} & m_{3}
    \end{pmatrix} = \delta_{m_{1}+m_{2}+m_{3},0} (-1)^{j_{1}-j_{2}-m_{3}} \\
    \cdot\sqrt{\frac{(j_{1}+j_{2}-j_{3})!(j_{1}-j_{2}+j_{3})!(-j_{1}+j_{2}+j_{3})!}{(j_{1}+j_{2}+j_{3}+1)!}} \\
    \cdot \sqrt{(j_{1}-m_{1})!(j_{1}+m_{1})!(j_{2}-m_{2})!(j_{2}+m_{2})!(j_{3}-m_{3})!(j_{3}+m_{3})!} \\
    \cdot \sum_{k=K}^{N} \frac{(-1)^{k}}{k!(j_{1}+j_{2}-j_{3}-k)!(j_{1}-m_{1}-k)!(j_{2}+m_{2}-k)!(j_{3}-j_{2}+m_{1}+k)!(j_{3}-j_{1}-m_{2}+k)!}
  \end{multline}
  with $K=\max(0,j_{2}-j_{3}-m_{1},j_{1}-j_{3}+m_{2})$ and
  $N=\min(j_{1}+j_{2}-j_{3},j_{1}-m_{1},j_{2}+m_{2})$, that is, the
  summation is taken such that the argument of each factorial in the
  denominator is non-negative. Factorials of negative numbers are
  taken as zero.

  For $kA\rightarrow\infty$ we have the asymptotic result
  \begin{align}
    C_{\sigma\mu\nu}^{sn(3)}(kA) &= \ordo\left(\frac{1}{kA}\right) \quad \text{for $\mu\neq\pm1$} \\
    C_{\sigma1\nu}^{sn(3)}(kA) &= \frac{\sqrt{(2n+1)(2\nu+1)}}{2} \iu^{\nu-n-1} \frac{\eu^{\iu kA}}{kA} + \ordo\left(\frac{1}{kA}\right) \\
    C_{\sigma,-1,\nu}^{sn(3)}(kA) &= \frac{\sqrt{(2n+1)(2\nu+1)}}{2} \iu^{\nu-n-1} (-1)^{s+\sigma} \frac{\eu^{\iu kA}}{kA} + \ordo\left(\frac{1}{kA}\right)     
  \end{align}
  The transformation of expansion coefficients is given by (assuming
  $r'>|A|$)
  \begin{multline}
    \Ev(\rv) = \sum_{cs\mu n} Q_{s\mu n}^{(c)}\Fv_{s\mu n}^{(c)}(r,\theta,\phi) \\
    = \sum_{c}\sum_{s=1}^{2}\sum_{n=1}^{\infty}\sum_{\mu=-n}^{n} Q_{s\mu n}^{(c)}\sum_{\sigma=1}^{2}\sum_{\nu=|\mu|,\nu\neq0}^{\infty} C_{\sigma\mu\nu}^{sn(1)}(kA)\Fv_{\sigma\mu\nu}^{(c)}(r',\theta',\phi') \\
    = \sum_{c}\sum_{s=1}^{2}\sum_{n=1}^{\infty}\sum_{\mu=-\infty}^{\infty} \chi(|\mu|\leq n)Q_{s\mu n}^{(c)}\sum_{\sigma=1}^{2}\sum_{\nu=|\mu|,\nu\neq0}^{\infty} C_{\sigma\mu\nu}^{sn(1)}(kA)\Fv_{\sigma\mu\nu}^{(c)}(r',\theta',\phi') \\
    = \sum_{c}\sum_{\sigma=1}^{2}\sum_{\mu=-\infty}^{\infty}\sum_{\nu=|\mu|,\nu\neq0}^{\infty}\sum_{s=1}^{2}\sum_{n=1}^{\infty} \chi(|\mu|\leq n)Q_{s\mu n}^{(c)} C_{\sigma\mu\nu}^{sn(1)}(kA)\Fv_{\sigma\mu\nu}^{(c)}(r',\theta',\phi') \\
    = \sum_{c}\sum_{\sigma=1}^{2}\sum_{\mu=-\infty}^{\infty}\sum_{\nu=1}^{\infty}\sum_{s=1}^{2}\sum_{n=1}^{\infty} \chi(|\mu|\leq\nu)\chi(|\mu|\leq n)Q_{s\mu n}^{(c)} C_{\sigma\mu\nu}^{sn(1)}(kA)\Fv_{\sigma\mu\nu}^{(c)}(r',\theta',\phi') \\
    = \sum_{c}\sum_{\sigma=1}^{2}\sum_{\nu=1}^{\infty}\sum_{\mu=-\nu}^{\nu}\sum_{s=1}^{2}\sum_{n=|\mu|,n\neq0}^{\infty} Q_{s\mu n}^{(c)} C_{\sigma\mu\nu}^{sn(1)}(kA)\Fv_{\sigma\mu\nu}^{(c)}(r',\theta',\phi') \\
    = \sum_{c\sigma\mu\nu} \underbrace{\left[\sum_{s=1}^{2}\sum_{n=|\mu|,n\neq0}^{\infty}Q_{s\mu n}^{(c)}C_{\sigma\mu\nu}^{sn(1)}(kA)\right]}_{=Q_{\sigma\mu\nu}^{(c)'}} \Fv_{\sigma\mu\nu}^{(c)}(r',\theta',\phi')
  \end{multline}
  where the characteristic function $\chi(|\mu|\leq n)=1$ when
  $|\mu|\leq n$ and 0 otherwise. In practice the sums of $n$ and $\nu$
  are finite with $n\leq N$ and $\nu\leq N'$, where we estimate $N'$ as
  the smallest integer larger than $N + kA$. For $r'<|A|$ we have
  \begin{multline}
    \Ev(\rv) = \sum_{cs\mu n} Q_{s\mu n}^{(c)}\Fv_{s\mu n}^{(c)}(r,\theta,\phi) \\
    = \sum_{\sigma\mu\nu} \underbrace{\left[\sum_{c}\sum_{s=1}^{2}\sum_{n=|\mu|,n\neq0}^{\infty}Q_{s\mu n}^{(c)}C_{\sigma\mu\nu}^{sn(c)}(kA)\right]}_{=Q_{\sigma\mu\nu}^{(1)'}} \Fv_{\sigma\mu\nu}^{(1)}(r',\theta',\phi')
  \end{multline}

\item Scattering matrix of $P$-port antenna:
  \begin{equation}
    \Ev(r,\theta,\phi) = \frac{k}{\sqrt{\eta}} \sum_{j=1}^{J} a_{j}\Fv_{j}^{(4)}(r,\theta,\phi) + b_{j}\Fv_{j}^{(3)}(r,\theta,\phi)
  \end{equation}
  \begin{equation}
    \begin{pmatrix}
      \boldsymbol{\Gamma} & \mat{R} \\
      \mat{T} & \mat{S} 
    \end{pmatrix}
    \begin{pmatrix}
      \vec{v} \\
      \vec{a}
    \end{pmatrix} =
    \begin{pmatrix}
      \vec{w} \\
      \vec{b}
    \end{pmatrix}
  \end{equation}
  Source scattering (T-matrix)
  \begin{equation}
    \Ev(r,\theta,\phi) = \frac{k}{\sqrt{\eta}} \sum_{j=1}^{J} a_{j}'\Fv_{j}^{(1)}(r,\theta,\phi) + b_{j}'\Fv_{j}^{(3)}(r,\theta,\phi)
  \end{equation}
  \begin{equation}
    \begin{pmatrix}
      \boldsymbol{\Gamma}' & \mat{R}' \\
      \mat{T}' & \mat{S}' 
    \end{pmatrix}
    \begin{pmatrix}
      \vec{v} \\
      \vec{a}'
    \end{pmatrix} =
    \begin{pmatrix}
      \vec{w} \\
      \vec{b}'
    \end{pmatrix} \quad\Rightarrow\quad
    \begin{pmatrix}
      \boldsymbol{\Gamma}' & \mat{R}' \\
      \mat{T}' & \mat{S}' 
    \end{pmatrix} = 
    \begin{pmatrix}
      \boldsymbol{\Gamma} & \frac{1}{2}\mat{R} \\
      \mat{T} & \frac{1}{2}(\mat{S}-\mat{I}) 
    \end{pmatrix}
  \end{equation}
\end{itemize}

\section{Definitions used in SciPy}

\begin{itemize}
\item Associated Legendre polynomials (scipy.special.lpmv)
  \begin{equation}
    \texttt{lpmv(m, n, x)} = \operatorname{P}_{n}^{m}(x) = (-1)^{m}(1 - x^{2})^{m/2} \frac{\diff^{m}P_{n}(x)}{\diff x^{m}}
  \end{equation}
  where the Legendre polynomials $P_{n}(x)$ are given by Rodrigues'
  formula (same as in Hansen)
  \begin{equation}
    P_{n}(x) = \frac{1}{2^{n}n!}\frac{\diff^{n}}{\diff x^{n}}(x^{2} - 1)^{n}
  \end{equation}
  This definition of $\operatorname{P}_{n}^{m}(x)$ is the same as in
  \cite[8.6.6 and 8.6.18]{Abramowitz+Stegun1970} and differs from the
  definition of $P_{n}^{m}(x)$ in Hansen by the Condon-Shortley phase
  function $(-1)^{m}$,
  $\operatorname{P}_{n}^{m}(x) = (-1)^{m}P_{n}^{m}(x)$.

  The derivative can be computed from the recursion relation
  \cite[8.5.2]{Abramowitz+Stegun1970}
  \begin{equation}
    (z^{2}-1)\frac{\diff\operatorname{P}_{\nu}^{\mu}(z)}{\diff z} = (\nu+\mu)(\nu-\mu+1)(z^{2}-1)^{1/2}\operatorname{P}_{\nu}^{\mu-1}(z) - \mu z\operatorname{P}_{\nu}^{\mu}(z)
  \end{equation}
  Since we want to evaluate this on the cut $-1<x<1$, we follow
  \cite[8.3.2 and 8.3.4]{Abramowitz+Stegun1970} and identify
  $\operatorname{P}_{\nu}^{\mu}(x) =
  \eu^{\pm\iu\mu\pi/2}\operatorname{P}_{\nu}^{\mu}(x\pm\iu0)$ and
  $(z^{2}-1) = (1-(x\pm\iu0)^{2})\eu^{\pm\iu\pi} =
  (1-x^{2})\eu^{\pm\iu\pi}$. After multiplication with
  $\eu^{\pm\iu\mu\pi/2}$ the left hand side is then
  \begin{equation}
    (1-x^{2})\eu^{\pm\iu\pi}\eu^{\pm\iu\mu\pi/2}\frac{\diff\operatorname{P}_{\nu}^{\mu}(x\pm\iu0)}{\diff x} = -(1-x^{2})\frac{\diff\operatorname{P}_{\nu}^{\mu}(x)}{\diff x}
  \end{equation}
  and the right hand side is
  \begin{multline}
    (\nu+\mu)(\nu-\mu+1)(1-x^{2})^{1/2}\eu^{\pm\iu\pi/2}\eu^{\pm\iu\mu\pi/2}\operatorname{P}_{\nu}^{\mu-1}(x\pm\iu0) - \mu x\eu^{\pm\iu\mu\pi/2}\operatorname{P}_{\nu}^{\mu}(x\pm\iu0) \\
    = (\nu+\mu)(\nu-\mu+1)(1-x^{2})^{1/2}\eu^{\pm\iu\pi/2}\eu^{\pm\iu\pi/2}\eu^{\pm\iu(\mu-1)\pi/2}\operatorname{P}_{\nu}^{\mu-1}(x\pm\iu0) - \mu x\eu^{\pm\iu\mu\pi/2}\operatorname{P}_{\nu}^{\mu}(x\pm\iu0) \\
    = - (\nu+\mu)(\nu-\mu+1)(1-x^{2})^{1/2}\operatorname{P}_{\nu}^{\mu-1}(x) - \mu x\operatorname{P}_{\nu}^{\mu}(x) 
  \end{multline}
  Dividing left and right hand side by $-(1-x^{2})$ and restricting
  our attention to integer order $\mu=m$ and degree $\nu=n$, we find
  \begin{equation}
    \frac{\diff\operatorname{P}_{n}^{m}(x)}{\diff x} = \frac{(n+m)(n-m+1)}{(1-x^{2})^{1/2}}\operatorname{P}_{n}^{m-1}(x) + \frac{mx}{1-x^{2}}\operatorname{P}_{n}^{m}(x)
  \end{equation}
  This relation can be used to compute all derivatives for positive
  integers $1\leq m\leq n$, $n\geq1$. For $m=0$, we need to consider
  \cite[8.2.5]{Abramowitz+Stegun1970}
  \begin{equation}
    \operatorname{P}_{\nu}^{-\mu}(z) = \frac{\Gamma(\nu-\mu+1)}{\Gamma(\nu+\mu+1)}\left[ \operatorname{P}_{\nu}^{\mu}(z) - \frac{2}{\pi}\eu^{-\iu\mu\pi}\sin(\mu\pi)Q_{\nu}^{\mu}(z)\right]
  \end{equation}
  For integers $\nu=n$ and $\mu=m$, we have
  $\Gamma(n-m+1)/\Gamma(n+m+1) = (n-m)!/(n+m)!$ and $\sin(m\pi)=0$,
  hence for $-1<x<1$ we have
  $\eu^{\pm\iu(-m)\pi/2}\operatorname{P}_{n}^{-m}(x) =
  \frac{(n-m)!}{(n+m)!}\eu^{\pm\iu
    m\pi/2}\operatorname{P}_{n}^{m}(x)$, and in particular
  $\operatorname{P}_{n}^{-1}(x) =
  -\frac{1}{n(n+1)}\operatorname{P}_{n}^{1}(x)$. Hence, the derivative
  can be written
  \begin{equation}
    \frac{\diff\operatorname{P}_{n}^{m}(x)}{\diff x} =
    \begin{cases}
      -\frac{1}{\sqrt{1-x^{2}}}\operatorname{P}_{n}^{1}(x) & m=0 \\
      \frac{(n+m)(n-m+1)}{\sqrt{1-x^{2}}}\operatorname{P}_{n}^{m-1}(x) + \frac{mx}{1-x^{2}}\operatorname{P}_{n}^{m}(x) & m=1,2,\ldots
    \end{cases}
    \label{eq:dPmndx}
  \end{equation}  
  To summarize, the functions $P_{n}^{m}(x)$ used in Hansen correspond
  to the functions $\operatorname{P}_{n}^{m}(x)$ implemented in
  scipy.special.lpmv as
  \begin{align}
    P_{n}^{m}(x) &= (-1)^{m}\operatorname{P}_{n}^{m}(x) \\
    \frac{\diff P_{n}^{m}(\cos\theta)}{\diff\theta} &= (-\sin\theta)\frac{\diff P_{n}^{m}(\cos\theta)}{\diff(\cos\theta)} = (-1)^{m}(-\sin\theta)\frac{\diff\operatorname{P}_{n}^{m}(\cos\theta)}{\diff(\cos\theta)}
  \end{align}
  with $\frac{\diff\operatorname{P}_{n}^{m}(x)}{\diff x}$ defined in
  \eqref{eq:dPmndx}.

\item Some special functions can be sensitive to high orders depending
  on implementation. It has proven good to use standard libraries
  rather than implementing text book formulas, for instance
  \verb+wigners+ to implement the Wigner 3-j symbols, and the
  \verb+scipy.special.jacobi+ for implementing the rotation
  coefficients $d^{n}_{\mu m}(\theta)$.
\end{itemize}

\section{Implementation of spherical vector wave library}

\begin{itemize}
\item Handling of apparent singularity at $\theta=0$ and $\theta=\pi$:
  identify arguments with $\theta=0$ and $\theta=\pi$, change to
  $\theta=\delta$ and $\theta=\pi-\delta$ for small $\delta$.

\item If we know the tangential components of the outgoing electric
  field on a spherical surface, we can use the orthogonality of
  tangential components \eqref{eq:orthotang} to compute
  \begin{multline}
    \int_{\phi=0}^{2\pi}\int_{\theta=0}^{\pi} \Big[ \{\Ev(r,\theta,\phi)\cdot\thetauv\} \{\Fv_{s,-m,n}^{(c)}(r,\theta,\phi)\cdot\thetauv\} \\
    + \{\Ev(r,\theta,\phi)\cdot\phiuv\} \{\Fv_{s,-m,n}^{(c)}(r,\theta,\phi)\cdot\phiuv\} \Big] \sin\theta\diff\theta\diff\phi \\
    = (-1)^{m}R_{sn}^{(3)}(kr)R_{sn}^{(c)}(kr)Q_{smn}^{(3)}
  \end{multline}
  The integral for the $\theta$ component can be written
  \begin{multline}
    \int_{\phi=0}^{2\pi}\int_{\theta=0}^{\pi} \{\Ev(r,\theta,\phi)\cdot\thetauv\} \{\Fv_{s,-m,n}^{(c)}(r,\theta,\phi)\cdot\thetauv\} \sin\theta\diff\theta\diff\phi \\
    = \int_{\theta=0}^{\pi} \left[ \int_{\phi=0}^{2\pi}\{\Ev(r,\theta,\phi)\cdot\thetauv\} \eu^{-\iu m\phi}\diff\phi\right] \{\Fv_{s,-m,n}^{(c)}(r,\theta,0)\cdot\thetauv\} \sin\theta\diff\theta
  \end{multline}
  and similarly for the $\phi$ component. This illustrates that the
  $\phi$ integral can be computed from the azimuth Fourier transform
  of the field.

  The far field is
  \begin{equation}
    \Fv(\theta,\phi) = \frac{1}{\sqrt{4\pi\eta}} \sum_{smn}Q_{smn}^{(3)}\Kv_{smn}(\theta,\phi) = \frac{1}{\sqrt{4\pi\eta}}\sum_{smn}Q_{smn}^{(3)}\Kv_{smn}(\theta,0)\eu^{\iu m\phi}
  \end{equation}
  which can be computed on a $\theta$-$\phi$ grid by an outer product
  between vectors $\Kv_{smn}(\theta,0)$ and $\eu^{\iu m\phi}$.
  
\item The \texttt{scipy.special.lpmv} implementation overflows for
  large orders, $N\geq86$. The package \texttt{pyshtools} implements
  normalized Legendre polynomials to high orders (2800 claimed),
  however it seems slower than \texttt{scipy.special.lpmv}.
\end{itemize}

\section{Application example: probe correction}

The received signal is \cite[Eq. (3.10)]{Hansen1988}
\begin{equation}
  w(A,\chi,\theta,\phi) = \frac{v}{2}\sum_{smn,\sigma\mu\nu} T_{smn}\eu^{\iu m\phi} d_{\mu m}^{n}(\theta) \eu^{\iu\mu\chi} C_{\sigma\mu\nu}^{sn(3)}(kA) R_{\sigma\mu\nu}^{\mrm{p}}
  \label{eq:transmissionformula}
\end{equation}
The purpose is to determine the transmitting coefficients $T_{smn}$ of
the antenna under test (AUT) given $w$ and the receiving coefficients
of the probe, $R_{\sigma\mu\nu}^{\mrm{p}}$. Assuming we have
measurements for all angles $0<\chi<2\pi$, $0<\theta<\pi$, and
$0<\phi<2\pi$, we can use the orthogonalities of the exponential
function and the rotation coefficients $d_{\mu m}^{n}(\theta)$ to find
\cite[Eqns. (4.51), (4.52)]{Hansen1988}
\begin{multline}
  w_{\mu m}^{n} = \frac{2n+1}{2(2\pi)^{2}}\int_{\chi=0}^{2\pi}\int_{\theta=0}^{\pi}\int_{\phi=0}^{2\pi} w(A,\chi,\theta,\phi) \eu^{-\iu\mu\chi}d_{\mu m}^{n}(\theta)\eu^{-\iu m\phi} \sin\theta\diff\phi\diff\theta\diff\chi \\
  = \sum_{s} vT_{smn}P_{s\mu n}(kA) = vT_{1mn}P_{1\mu n}(kA) + vT_{2mn}P_{2\mu n}(kA)
\end{multline}
where we introduced the probe response constants as
\begin{equation}
  P_{s\mu n}(kA) = \frac{1}{2}\sum_{\sigma\nu}C_{\sigma\mu\nu}^{sn(3)}(kA)R_{\sigma\mu\nu}^{\mrm{p}}
\end{equation}
It is common that the probe has nonzero coefficients
$R_{\sigma\mu\nu}^{\mrm{p}}$ for $\mu=\pm1$ only. In this case the
equations reduce to
\begin{align}
  vT_{1mn}P_{11n}(kA) + vT_{2mn}P_{21n}(kA) &= w_{1m}^{n}(A) \\
  vT_{1mn}P_{1,-1,n}(kA) + vT_{2mn}P_{2,-1,n}(kA) &= w_{-1,m}^{n}(A)
\end{align}
which can be solved for $T_{smn}$. For a $\mu=\pm1$ probe, we do not
need to measure all angles $0<\chi<2\pi$, since we then have
$w(A,\chi,\theta,\phi) = w_{1}\eu^{\iu\chi} + w_{-1}\eu^{-\iu\chi}$,
where \cite[p. 110]{Hansen1988}
\begin{align}
  w_{1}(A,\theta,\phi) = \frac{1}{2} \left\{ w(A,0,\theta,\phi) - \iu w(A,\pi/2,\theta,\phi) \right\} \\
  w_{-1}(A,\theta,\phi) = \frac{1}{2} \left\{ w(A,0,\theta,\phi) + \iu w(A,\pi/2,\theta,\phi) \right\} 
\end{align}
Hence, for a $\mu=\pm1$ probe it is sufficient to measure only for
$\chi=0$ and $\chi=\pi/2$.

We now turn to the formulation of the equations assuming we are
restricted to the angle $\chi$ taking only two values, 0 or $\pi/2$ (or
possibly $-\pi/2$), corresponding to the two orientations of the probe
which determines polarization of the measured signal $w$.
\begin{equation}
  w_{m}(A,\chi,\theta) = \frac{1}{2\pi}\int_{\phi=0}^{2\pi} \eu^{-\iu m\phi}w(A,\chi,\theta,\phi) \diff\phi
\end{equation}
we rewrite \eqref{eq:transmissionformula} as
\begin{equation}
  w_{m}(A,\chi,\theta,\phi) = \sum_{sn,\sigma\mu\nu} vT_{smn} d_{\mu m}^{n}(\theta) \eu^{\iu\mu\chi} P_{s\mu n}(kA)
\end{equation}
Multiplying by $d_{\mu'm}^{n'}(\theta)\sin\theta$ and integrating, we find
\begin{multline}
  \int_{\theta=0}^{\pi}w_{m}(A,\chi,\theta) d_{\mu'm}^{n'}(\theta) \sin\theta\diff\theta\diff\phi \\
  = \sum_{sn,\sigma\mu\nu} vT_{smn}\eu^{\iu\mu\chi} P_{s\mu n}(kA)\int_{\theta=0}^{\pi}  d_{\mu m}^{n}(\theta) d_{\mu'm}^{n'}(\theta) \sin\theta \diff\theta 
\end{multline}
The real factors
\begin{equation}
  D_{\mu\mu'm}^{nn'} = \int_{\theta=0}^{\pi}  d_{\mu m}^{n}(\theta) d_{\mu'm}^{n'}(\theta) \sin\theta \diff\theta 
\end{equation}
can be precomputed and stored. From \cite[Eq. (A2.10)]{Hansen1988}, it
is clear that $D_{\mu\mu m}^{nn'} = \frac{2}{2n+1}\delta_{nn'}$. We
then have the equation system
\begin{align}
  \int_{\theta=0}^{\pi}w_{m}(A,0,\theta)d_{\mu'm}^{n'}(\theta)\sin\theta\diff\theta &= \sum_{sn} vT_{smn} \sum_{\sigma\mu\nu} P_{s\mu n}(kA) D_{\mu\mu'n}^{nn'} \label{eq:probecorrectionA} \\
  \int_{\theta=0}^{\pi}w_{m}(A,\pi/2,\theta)d_{\mu'm}^{n'}(\theta)\sin\theta\diff\theta &= \sum_{sn} vT_{smn} \sum_{\sigma\mu\nu} \eu^{\iu\mu\pi/2} P_{s\mu n}(kA) D_{\mu\mu'n}^{nn'} \label{eq:probecorrectionB}
\end{align}
To determine $R_{\sigma\mu\nu}^{\mrm{p}}$, we first assume the probe
is reciprocal and exists in two identical realizations. We then have \cite[Eq. (3.21)]{Hansen1988}
\begin{equation}
  T_{\sigma,-\mu,\nu}^{\mrm{p}} = (-1)^{\mu} R_{\sigma\mu\nu}^{\mrm{p}}
\end{equation}
However, the probe used as AUT is also rotated by $\pi$ around the
$y$-axis so that it is pointing towards the probe. Since
\cite[Eq. (A2.16)]{Hansen1988}
\begin{equation}
  d_{\mu m}^{n}(\pi) = (-1)^{n+m}\delta_{\mu,-m}
\end{equation}
we then have
\begin{equation}
  T_{s\mu n} = \sum_{m=-n}^{n}T_{smn}^{\mrm{p}}d_{\mu m}^{n}(\pi) = (-1)^{n-\mu}T_{s,-\mu,n}^{\mrm{p}} = (-1)^{n-\mu}(-1)^{\mu}R_{s\mu n}^{\mrm{p}} = (-1)^{n} R_{s\mu n}^{\mrm{p}}
\end{equation}
To determine the probe coefficients $R_{smn}^{\mrm{p}}$, now make
measurements $w^{\mrm{p}}(A,\chi,\theta,\phi)$ using the probe
antennas both as probe and AUT. We can then start with a guess for the
probe receiving coefficients as a Hertzian dipole, that is,
\begin{equation}
  R_{211}^{\mrm{p},0} = -R_{2,-1,1}^{\mrm{p},0} = -1/\sqrt{2}, \quad \text{all other coefficients zero}
\end{equation}
Use equations \eqref{eq:probecorrectionA} and
\eqref{eq:probecorrectionB} to solve for $T_{smn}^{1}$, and compute
$R_{smn}^{\mrm{p},1} = (-1)^{n}T_{smn}^{1}$. Use $R_{smn}^{\mrm{p},1}$
to solve for $T_{smn}^{2}$, and iterate until convergence.

As it turns out, there seems to be a tendency for the iteration to get
stuck in a bifurcation, where
$R_{smn}^{\mrm{p},i}=R_{smn}^{\mrm{p,e}}$ for $i$ even and
$R_{smn}^{\mrm{p},i}=R_{smn}^{\mrm{p,o}}$ for $i$ odd. In order to
counteract this, the iteration can be slightly changed so that once
$R_{smn}^{\mrm{p},i+1}$ has been computed, it is recomputed as
\begin{equation}
  R_{smn}^{\mrm{p},i+1} \leftarrow \frac{1}{2}\left( R_{smn}^{\mrm{p},i+1} + R_{smn}^{\mrm{p},i} \right)
\end{equation}

\bibliography{add,total}
\bibliographystyle{teorel}

\end{document}